\subsection{Bidirectional Data Flow Pipelines}

Generator-based coroutines extend \texttt{yield} into a bidirectional channel: the caller can inject values, exceptions, and shutdown signals into a suspended frame.

\subsubsection{Consumers and Transformers: Feeding In-Flight Data via the \texttt{.send()} Interface}

\texttt{gen.send(value)} resumes the generator and sets \texttt{yield} to receive \texttt{value}. The first \texttt{send} must use \texttt{None} (or precede any non-\texttt{None} send with \texttt{next()}), because execution must reach the first \texttt{yield} before a value can be injected.

\begin{codebox}[]{python}{Consumers and Transformers: Feeding In-Flight Data via the .send() Interface}{bid01}
def echo():
    while True:
        received = yield
        print(f"got: {received!r}")

gen = echo()
next(gen)
gen.send("D'oh!")
gen.send("Ni!")
\end{codebox}


\begin{iobox}{Expected Output}
\texttt{got: "D'oh!"} \\
\texttt{got: 'Ni!'} \\
\end{iobox}




\subsubsection{Exception Injection with \texttt{.throw()} and Controlled Shutdown with \texttt{.close()}}

\texttt{gen.throw(exc)} injects an exception at the suspension point. \texttt{gen.close()} raises \texttt{GeneratorExit} inside the generator, triggering cleanup. Both operate on the same heap frame as ordinary resumption.

\subsubsection{Cooperative Multitasking: Voluntary Suspension Instead of Preemptive Scheduling}

Generators yield control explicitly. No OS thread preemption occurs; the caller decides when to resume. This cooperative model avoids locks on local state but requires disciplined scheduling by the driving loop.
